\documentclass[12pt,a4paper]{article}
\usepackage[UTF8]{ctex}
\usepackage{geometry}
\usepackage{graphicx}
\usepackage{booktabs}
\usepackage{longtable}
\usepackage{array}
\usepackage{multirow}
\usepackage{amsmath}
\usepackage{hyperref}
\usepackage{xcolor}
\usepackage{tikz}
\usepackage{enumitem}
\usepackage{fancyhdr}
\usepackage{titlesec}
\usepackage{tabularx}

\geometry{left=2.5cm, right=2.5cm, top=2.5cm, bottom=2.5cm}

\pagestyle{fancy}
\fancyhf{}
\fancyhead[L]{\small 软件项目开发与实践}
\fancyhead[R]{\small 讨论记录二}
\fancyfoot[C]{\thepage}
\renewcommand{\headrulewidth}{0.4pt}

\titleformat{\section}{\Large\bfseries\color{blue!70!black}}{第\,\thesection\,部分}{0.8em}{}
\titleformat{\subsection}{\large\bfseries}{}{0em}{}
\titleformat{\subsubsection}{\normalsize\bfseries}{}{0em}{}

\hypersetup{
    colorlinks=true,
    linkcolor=blue!70!black,
    urlcolor=blue!60!black,
}

\begin{document}

% ==================== 标题 ====================
\begin{center}
{\LARGE\bfseries 课程学习小组讨论记录}
\end{center}

\vspace{0.5cm}

\noindent\textbf{【讨论主题】}：AI 对软件开发的影响及工具分享

\vspace{0.3cm}
\noindent\textbf{【课程名称】}：软件项目开发与实践

\vspace{0.3cm}
\noindent\textbf{【讨论日期】}：2026年6月

\vspace{0.3cm}
\noindent\textbf{【讨论时长】}：2小时

\vspace{0.3cm}
\noindent\textbf{【主持人】}：谢宇鹏

\vspace{0.3cm}
\noindent\textbf{【参与人员】}：谢宇鹏（学号：5120243390）

\vspace{0.8cm}

% ==================== 一、讨论背景与目标 ====================
\section{讨论背景与目标}

\subsection{【讨论背景】}

最近两年，AI编程工具发展得特别快。最早的时候，AI只能帮你补全几个单词、几行代码，就像手机打字时的联想功能一样。但到了2024-2025年，Cursor、Claude Code、Trae这些新工具出来了，AI的能力有了质的飞跃——它不再只是帮你补代码了，而是能\textbf{理解你想做什么}，帮你规划怎么做、帮你写出完整的功能模块、帮你找Bug修Bug。

我们这门课用的是B模式（AI智能协作开发），就是说游戏项目要在AI的辅助下来完成。这不只是「让AI帮我写代码」那么简单——我们需要理解AI擅长什么、不擅长什么，学会怎么跟AI「合作」，让它成为我们的好帮手而不是替代品。

我们做的「星露谷战斗模块」项目，从需求讨论到代码实现，全程都有AI参与。这次讨论就是来总结：AI到底怎么影响了软件开发？我们用了哪些AI工具？怎么跟AI配合才能做出好东西？

\subsection{【讨论目标】}

\begin{enumerate}[leftmargin=2em]
    \item 分析AI对软件开发各个阶段（想需求、做设计、写代码、找Bug）带来了什么变化
    \item 对比我们用过的几款AI工具，说说各自的好坏
    \item 总结跟AI合作时的有效方法和技巧
    \item 讨论「怎么给AI下指令」这件事为什么重要
    \item 聊聊AI时代程序员的角色会怎么变
\end{enumerate}


% ==================== 二、主要讨论内容 ====================
\section{主要讨论内容}

\subsection{AI 怎么影响了软件开发的各个环节}

\subsubsection{想需求的阶段}

\begin{itemize}[leftmargin=2em]
    \item \textbf{好处}：你只需要跟AI说「我想做一个类似星露谷的打怪小游戏」，它就能帮你展开成一份结构化的需求清单——要有哪些功能、什么场景、哪些角色，帮你从一团模糊的想法变成清晰的列表
    \item \textbf{好处}：AI能帮你生成需求文档的框架，用专业的格式写出来，省去很多排版的功夫
    \item \textbf{要注意}：AI对「好不好玩」这种感性的东西理解不太到位，它更偏技术视角，所以游戏性方面还得我们自己判断
    \item \textbf{要注意}：AI列出来的需求可能太多太杂，得我们自己根据时间和能力来裁剪，不能什么都想要
    \item \textbf{我们怎么做的}：先跟AI讨论游戏定位（「做星露谷的战斗模块」），再让AI细化功能点（最后列出了14个功能需求），然后我们自己调整优先级和细节
\end{itemize}

\subsubsection{做设计的阶段}

\begin{itemize}[leftmargin=2em]
    \item \textbf{好处}：告诉AI你的技术栈和偏好，它能快速帮你画出模块怎么划分、类怎么设计、数据怎么流动
    \item \textbf{好处}：你可以让AI比较好几种不同的设计方案，帮你分析每种的优缺点
    \item \textbf{要注意}：AI有时候会「过度设计」——明明简单的东西它给你搞得很复杂，引入不必要的结构
    \item \textbf{要注意}：AI可能不了解你的具体限制（比如Pygame有什么坑、课程给的时间有多少），设计出来的东西可能不接地气
    \item \textbf{我们怎么做的}：告诉AI我们要用面向对象、分层架构，让它画出类结构图，然后我们自己审查，把不必要的删掉
\end{itemize}

\subsubsection{写代码的阶段}

\begin{itemize}[leftmargin=2em]
    \item \textbf{好处}：AI能快速生成那些模式化的代码（比如精灵类、状态机、碰撞检测的框架），省掉大量重复劳动
    \item \textbf{好处}：AI对Pygame这些库的API很熟，调用方式基本不会错
    \item \textbf{好处}：遇到Bug的时候，把报错信息发给AI，它能很快帮你定位问题、给出修复方案
    \item \textbf{要注意}：AI写的代码可能有逻辑错误，比如碰撞检测的边界情况没处理好
    \item \textbf{要注意}：AI的代码「能跑但不一定好」——有时候写得很臃肿，需要我们自己重构优化
    \item \textbf{我们怎么做的}：一个模块一个模块地来——先做地图，跑通了再做玩家，再做怪物，再做武器。每做完一个模块就立刻运行测试，发现问题马上反馈给AI修
\end{itemize}

\subsubsection{找Bug的阶段}

\begin{itemize}[leftmargin=2em]
    \item \textbf{好处}：AI能根据功能需求帮你列出测试清单，不用自己一个一个想
    \item \textbf{好处}：把错误堆栈信息发给AI，它能快速指出哪里出了问题，给出修复建议
    \item \textbf{要注意}：AI修了一个Bug可能又引入了新的Bug，修完之后得全面测一遍
    \item \textbf{我们怎么做的}：把运行时报错的完整信息复制给AI，AI给出修复方案，我们先理解再应用，不是盲目照搬
\end{itemize}

\subsection{我们用过的AI工具对比}

\begin{longtable}{p{2cm}p{1.8cm}p{3cm}p{3.5cm}p{3cm}}
\toprule
\textbf{工具} & \textbf{类型} & \textbf{我们怎么用的} & \textbf{好在哪} & \textbf{不好在哪} \\
\midrule
\endhead
Claude Code & 命令行工具 & 讨论需求、设计架构、写代码、写文档 & 理解长文本能力强，适合深入讨论；生成代码质量高；能直接帮你改文件 & 需要付费 \\
GitHub Copilot & 编程助手插件 & 日常写代码时自动补全 & 跟编辑器集成流畅，对各种库都很熟，补全速度快 & 只会补全，不会主动帮你规划 \\
Cursor & 智能编辑器 & 全流程开发、直接编辑文件 & 体验流畅，能直接创建和修改多个文件，智能程度高 & 有一定学习成本 \\
Trae & 国产智能编辑器 & 备选方案 & 免费使用，中文支持好，功能还在不断完善 & 免费版有次数限制 \\
ChatGPT & 网页对话 & 技术调研、解释原理 & 灵活，知识面广，什么都能聊 & 不实时，代码上下文有限 \\
\bottomrule
\end{longtable}

\textbf{我们的心得}：
\begin{itemize}[leftmargin=2em]
    \item \textbf{深入讨论和设计}：用Claude Code最合适，它能理解很长的上下文，跟你来回讨论
    \item \textbf{日常写代码}：GitHub Copilot或Cursor，写代码时自动补全，效率很高
    \item \textbf{从头到尾做一个项目}：Cursor或Claude Code的Agent模式，能直接帮你操作文件
    \item \textbf{预算有限}：Trae（免费）+ ChatGPT（免费版）也能用
\end{itemize}

\subsection{跟AI合作的方法和技巧}

\subsubsection{想需求的时候怎么跟AI聊}

\begin{enumerate}[leftmargin=2em]
    \item \textbf{从模糊到具体}：先说个大概（「做个类似星露谷的打怪游戏」），再一步步细化（「要有哪几种武器？几种怪物？几张地图？」）
    \item \textbf{让AI帮你展开}：跟AI说「帮我列一个完整的功能清单」，它会帮你生成结构化的列表
    \item \textbf{自己裁剪}：AI列出来的东西可能太多，根据你的实际能力和时间来决定先做什么后做什么
    \item \textbf{记下来}：把讨论结果保存成文档，后面写代码的时候可以参考
\end{enumerate}

\subsubsection{做设计的时候怎么跟AI聊}

\begin{enumerate}[leftmargin=2em]
    \item \textbf{说清楚限制}：告诉AI你用什么技术、有什么要求、不能用什么（「用Python + Pygame，面向对象，多文件结构，不用TMX地图格式」）
    \item \textbf{让AI解释}：不要只让AI给代码，让它解释为什么这么设计，你才能学到东西
    \item \textbf{让AI比较}：让AI列出几种方案，对比各自的优缺点，你再做决定
    \item \textbf{让AI出文档}：让AI帮你生成设计文档的框架，你再补充技术细节
\end{enumerate}

\subsubsection{写代码的时候怎么跟AI聊}

\begin{enumerate}[leftmargin=2em]
    \item \textbf{一个一个来}：不要让AI一次性把整个项目写出来，而是先做地图，再做玩家，再做怪物，一个模块一个模块地来
    \item \textbf{告诉AI上下文}：写新模块的时候，告诉AI已有的类和接口长什么样，确保新代码能跟旧代码对接
    \item \textbf{写完就测}：每做完一个模块立刻跑一遍，有问题马上反馈给AI修
    \item \textbf{保存版本}：每次AI生成的代码确认没问题后立刻提交版本记录，方便以后回退
\end{enumerate}

\subsubsection{找Bug的时候怎么跟AI聊}

\begin{enumerate}[leftmargin=2em]
    \item \textbf{给完整信息}：把完整的报错信息复制给AI，不要只说「报错了」
    \item \textbf{说清楚期望}：告诉AI「我期望它怎么做，但实际上它怎么了」，AI才能准确判断问题
    \item \textbf{先理解再改}：AI给出修复方案后，先搞懂为什么这么改，再应用，不要盲目复制粘贴
    \item \textbf{修完要验证}：修了一个Bug后，检查其他功能有没有被影响
\end{enumerate}

\subsection{怎么给AI下指令（提示词技巧）}

在B模式开发中，你怎么跟AI说话，直接决定了AI输出的质量。

\subsubsection{好的指令长什么样}

\begin{enumerate}[leftmargin=2em]
    \item \textbf{目标明确}：「请用Python实现一个怪物自动绕开障碍物追玩家的寻路功能」比「帮我写个追人功能」好得多
    \item \textbf{有约束条件}：「在20$\times$11的格子地图上，用两个格子之间的距离来估算远近」
    \item \textbf{提供背景}：「我正在做Pygame游戏，已经有一个Entity基类，里面有移动和碰撞检测方法」
    \item \textbf{要求格式}：「请返回完整的Python代码，加上中文注释」
    \item \textbf{迭代改进}：「这段代码能跑了，但太慢了，请优化，减少每帧的计算量」
\end{enumerate}

\subsubsection{常见的错误指令}

\begin{enumerate}[leftmargin=2em]
    \item \textbf{太模糊}：「帮我做个游戏」——AI不知道做什么类型、什么规模
    \item \textbf{没约束}：「写个AI」——没有指定用什么语言、什么框架、做什么功能
    \item \textbf{一次太多}：「帮我做整个项目」——应该拆成一个个小任务
    \item \textbf{不给反馈}：AI第一次输出不满意时，应该说哪里不满意，而不是换个方式重新问
\end{enumerate}

\subsection{AI 时代程序员的角色变化}

\subsubsection{从「写代码」到「看代码}

以前程序员的核心工作是自己一行一行写代码。现在AI几秒钟就能生成你几个小时才能写完的代码量，所以写代码本身不再是核心竞争力了。程序员的角色更像是「审查员」——判断AI写的东西好不好、有没有问题、符不符合需求。

\subsubsection{从「做功能」到「做设计}

AI很擅长做具体的事情（比如「实现一个碰撞检测」），但不太擅长想整体方案（比如「怎么划分模块才能让系统好扩展」）。所以架构设计能力反而变得更重要了——这是AI替代不了的。

\subsubsection{从「会语法」到「懂原理}

AI能写出语法正确的代码，但程序员需要理解底层原理（比如为什么要先处理水平碰撞再处理垂直碰撞），才能判断AI写的东西合不合理。不懂原理的话，AI写了个有隐蔽Bug的代码你也看不出来。

\subsubsection{给AI下指令成了新技能}

能准确描述需求、把大任务拆成小任务、引导AI输出高质量的结果，成了AI时代程序员的核心技能之一。这其实跟传统的需求分析能力一脉相承，但要求更高——你不仅要懂业务，还要懂AI擅长什么、不擅长什么。


% ==================== 三、总结与反思 ====================
\section{总结与反思}

\subsection{讨论总结}

通过这次深入讨论，我们得出以下核心结论：

\begin{enumerate}[leftmargin=2em]
    \item \textbf{AI是工具不是替代}：AI大幅提高了开发效率，但不能替代我们的思考和判断
    \item \textbf{需求能力更重要}：「做什么」比「怎么做」更有价值，想清楚要什么是最大的竞争力
    \item \textbf{基础是前提}：数据结构、算法、设计模式这些基础知识是驾驭AI工具的前提，没基础就判断不了AI写的好不好
    \item \textbf{沟通是关键}：你怎么跟AI说话，决定了AI输出的质量——好的指令事半功倍
    \item \textbf{审查是底线}：所有AI生成的代码都必须人工检查过才能用，不能盲目信任
\end{enumerate}

\subsection{对软件工程专业学生的启示}

\begin{enumerate}[leftmargin=2em]
    \item \textbf{打好基础}：数据结构、算法、操作系统、网络这些基础知识是你的核心壁垒，AI暂时替代不了
    \item \textbf{培养设计思维}：从「实现功能」转向「设计系统」，理解怎么划分模块、怎么管理依赖、怎么让系统好扩展
    \item \textbf{学会给AI下指令}：准确描述需求、拆解任务、引导AI输出好结果，这是新必备技能
    \item \textbf{建立审查能力}：能判断AI写的东西好不好，发现问题并修正
    \item \textbf{拥抱变化}：AI不会取代程序员，但会取代不会用AI的程序员
    \item \textbf{持续学习}：AI工具在快速迭代，保持学习和适应能力是长期竞争力
\end{enumerate}

\end{document}
